\documentclass[12pt]{article}
\usepackage{amsmath}
\usepackage{amssymb}
\usepackage{geometry}
\usepackage{enumitem}

\geometry{margin=1in}

\title{Hierarchical Linear Regression Model}
\author{}
\date{\today}

\begin{document}

\maketitle

\section{Model Overview}

This document describes a hierarchical linear regression model with individual-specific parameters, designed to model longitudinal data with age-dependent responses and individual covariates.

\section{Data Structure}

\begin{itemize}
    \item $N$: Total number of observations
    \item $K$: Number of individuals
    \item $y_{i,k}$: Response variable for observation $i$ of individual $k$
    \item $\text{age}_{i,k}$: Age at observation $i$ for individual $k$
    \item $\text{edu}_k$: Education level for individual $k$
    \item $\text{sex}_k$: Sex indicator for individual $k$
    \item $\text{cag}_k$: CAG repeat length for individual $k$
    \item $s_k$: Number of observations for individual $k$
\end{itemize}

\section{Statistical Model}

\subsection{Likelihood}

For individual $k = 1, \ldots, K$ and observations $i$ within individual $k$:

\begin{equation}
y_{i,k} \sim \mathcal{N}(\mu_{i,k}, \sigma^2)
\end{equation}

where the mean is structured as:

\begin{equation}
\mu_{i,k} = \beta_{\text{edu}} \cdot \text{edu}_k + \beta_{\text{sex}} \cdot \text{sex}_k + \beta_{\text{cag}} \cdot \text{cag}_k + \beta_k \cdot (\text{age}_{i,k} - \alpha_k)
\end{equation}

\subsection{Individual-specific Parameters}

Each individual $k$ has their own intercept and slope parameters:

\begin{align}
\alpha_k &= \exp(\alpha_{\text{raw},k}) \quad \text{(individual-specific age offset)} \\
\beta_k &= \exp(\beta_{\text{raw},k}) \quad \text{(individual-specific age slope)}
\end{align}

The exponential transformation ensures that $\alpha_k, \beta_k > 0$.

\subsection{Prior Distributions}

\subsubsection{Individual-specific Parameters}
\begin{align}
\alpha_{\text{raw},k} &\sim \mathcal{N}(4, 0.7^2) \\
\beta_{\text{raw},k} &\sim \mathcal{N}(-4, 0.3^2)
\end{align}

\subsubsection{Population-level Parameters}
\begin{align}
\sigma &\sim \text{Gamma}(2, 1) \\
\beta_{\text{edu}} &\sim \mathcal{N}(0, 10^2) \\
\beta_{\text{sex}} &\sim \mathcal{N}(0, 10^2) \\
\beta_{\text{cag}} &\sim \mathcal{N}(0, 10^2)
\end{align}

\section{Model Interpretation}

\subsection{Model Type}
This is a \textbf{hierarchical nonlinear growth curve model} with the following characteristics:

\begin{itemize}[leftmargin=*]
    \item \textbf{Hierarchical structure}: Each individual has their own growth parameters
    \item \textbf{Nonlinear age effect}: The age relationship is shifted by individual-specific $\alpha_k$
    \item \textbf{Individual covariates}: Education, sex, and CAG repeat length affect baseline levels
    \item \textbf{Positive constraints}: Log-normal parameterization ensures $\alpha_k, \beta_k > 0$
\end{itemize}

\subsection{Parameter Interpretation}

\begin{itemize}[leftmargin=*]
    \item $\alpha_k$: Individual-specific age offset or inflection point (when effects begin)
    \item $\beta_k$: Individual-specific rate of change with respect to age
    \item $\beta_{\text{edu}}$: Effect of education on baseline response
    \item $\beta_{\text{sex}}$: Effect of sex on baseline response  
    \item $\beta_{\text{cag}}$: Effect of CAG repeat length on baseline response
    \item $\sigma$: Residual standard deviation
\end{itemize}

\subsection{Potential Applications}

Given the variables present (age, education, sex, CAG repeats), this model is likely designed for:

\begin{itemize}[leftmargin=*]
    \item \textbf{Disease progression modeling}: Particularly neurodegenerative diseases
    \item \textbf{Huntington's disease research}: CAG repeats are strongly associated with Huntington's
    \item \textbf{Cognitive decline studies}: Individual trajectories of cognitive function over age
    \item \textbf{Biomarker progression}: Tracking disease biomarkers over time
\end{itemize}

\section{Model Advantages}

\begin{enumerate}[leftmargin=*]
    \item \textbf{Individual heterogeneity}: Accounts for individual differences in disease onset and progression
    \item \textbf{Covariate adjustment}: Controls for demographic and genetic factors
    \item \textbf{Flexible age relationships}: Allows for individual-specific age trajectories
    \item \textbf{Bayesian framework}: Incorporates prior knowledge and quantifies uncertainty
    \item \textbf{Hierarchical pooling}: Borrows strength across individuals for better parameter estimation
\end{enumerate}

\end{document}